\documentclass[11pt]{article}
\usepackage{fontspec}
\defaultfontfeatures{Ligatures=TeX}
\usepackage{microtype}
\usepackage[margin=1in]{geometry}
\usepackage{amsmath,amssymb}
\usepackage{graphicx}
\usepackage{booktabs,longtable}
\usepackage{array}
\usepackage{siunitx}
\usepackage[numbers,sort&compress]{natbib}
\usepackage[unicode,hidelinks]{hyperref}
\usepackage[capitalise,noabbrev]{cleveref}
\usepackage{xcolor}
\usepackage{url}
\graphicspath{{figures/}}
\setlength{\emergencystretch}{3em}
\providecommand{\tightlist}{\setlength{\itemsep}{0pt}\setlength{\parskip}{0pt}}
\hypersetup{
  pdftitle={Optical Diffractive Neural Networks: A Comprehensive Review of Architectures, Training, and Future Directions},
  pdfauthor={OptoMind}
}

\title{Optical Diffractive Neural Networks: A Comprehensive Review of Architectures, Training, and Future Directions}
\author{OptoMind}
\date{}

\begin{document}
\maketitle

\begin{abstract}
Optical diffractive neural networks transform computational imaging by replacing iterative electronic inference with passive light propagation through learned phase patterns; however, their transition from laboratory prototypes to deployable optical AI hinges on resolving the tension between physical realizability constraints (wavelength sensitivity, fabrication tolerance, static vs. reconfigurable hardware) and the algorithmic flexibility that originally motivated diffractive deep learning. The review argues that ODNNs represent a fundamentally different computing paradigm---not merely an optical implementation of digital neural networks---where the physics of diffraction imposes both unique capabilities (ultra-low-latency, energy-efficient inference at the speed of light) and hard constraints (wavelength dependence, fixed or slowly tunable phase profiles, limited nonlinearities). The field has evolved from single-function diffractive lenses to general-purpose trainable optical processors, but progress is bottlenecked by hardware limitations rather than algorithmic ones. A pathway forward requires hybrid electro-optical co-design, standardized benchmarking protocols, and fabrication-aware training strategies that treat physical constraints as features rather than afterthoughts. This review maps the architectural landscape, evaluates training paradigms against physical realizability, catalogs applications by performance maturity, and identifies the technical bottlenecks that must be resolved for ODNNs to achieve practical impact.
\end{abstract}

\noindent\textbf{Keywords:} diffractive deep neural network, optical diffractive neural network, all-optical neural network, diffractive optical element machine learning, phase modulation neural network, multi-layer diffractive optical network, training diffractive optical networks, gradient descent optical neural network

\section{Introduction}\label{introduction}

Driven by the relentless demand for higher computational throughput and lower energy consumption in artificial intelligence, intense interest has surged in alternative computing paradigms that transcend the limitations of conventional electronic processors. Among these emerging technologies, Optical Diffractive Neural Networks (ODNNs), also known as Diffractive Deep Neural Networks (D2NNs), have emerged as a distinct class of physical computers. Unlike traditional optical accelerators that merely speed up specific matrix operations within an electronic workflow, ODNNs fundamentally replace electronic weight matrices with spatially varying phase profiles. These systems implement linear transformations via Fresnel or Fraunhofer diffraction, creating a computing architecture where the laws of wave physics directly dictate computational expressivity. The foundational promise of the resulting design was demonstrated by early all-optical implementations using passive, 3D-printed diffractive layers, which showed that complex functions such as image classification could be executed at the speed of light without active power consumption during inference \citep{doi_10_1126_science_aat8084_e055772d}. This review systematically examines the evolution of ODNNs from theoretical concepts to functional prototypes, evaluating their potential to redefine the boundaries of optical artificial intelligence.

Optical Deep Neural Networks (ODNNs) are significant because they address the von Neumann bottleneck, which involves latency and energy overhead from moving data between memory and processing units in electronic systems. By encoding information in the phase and amplitude of an optical wavefront and processing it through free-space propagation, ODNNs offer massive parallelism and ultra-low latency. Recent theoretical advancements have further expanded this potential, demonstrating that large-scale nonlinear computation can be approximated using linear optics through optimized diffractive processors, effectively serving as universal function approximators for bandlimited nonlinear functions \citep{doi_10_1186_s43593_025_00113_w_ca9e8e1e}. Such capabilities suggest a pathway toward ultrafast, energy-efficient inference engines for applications ranging from real-time object recognition to all-optical encryption and compressive sensing. In this regime, the computational cost is largely decoupled from the complexity of the operation, governed instead by the physical propagation time of light through the device.

However, the transition from idealized simulations to robust physical deployment is fraught with significant challenges that temper these theoretical advantages. A central tension in ODNN research is the ``nonlinearity bottleneck'': while passive diffractive layers excel at linear transformations, they lack the intrinsic nonlinearity required for deep learning architectures unless supplemented by external mechanisms or hybrid electronic interfaces. Furthermore, the fidelity of ODNNs is critically dependent on the precision of their physical realization. The ``simulation-to-reality'' gap remains a pervasive obstacle, where performance predicted by scalar diffraction models degrades in physical prototypes due to fabrication tolerances, surface roughness, and material dispersion. Bridging this gap requires sophisticated training paradigms, such as hardware-in-the-loop optimization, where the physical device is included in the training loop to account for real-world imperfections \citep{doi_10_1117_12_3013482_d44587c1}. Additionally, the static nature of many fabricated platforms---such as 3D-printed polymers or lithographically defined metasurfaces---creates ``write-once'' hardware that lacks the reconfigurability essential for dynamic computing environments. While reconfigurable platforms based on liquid crystal spatial light modulators (SLMs) offer adaptability, they often introduce trade-offs in switching speed, power consumption, and diffraction efficiency \citep{doi_10_1002_adom_202300561_c4d7a5b0}.

This review provides a comprehensive scholarly analysis of the ODNN landscape, structured to guide the reader from fundamental physical principles to future research trajectories. We begin by establishing the physical architecture of diffractive neural layers, explicitly framing the interplay between diffraction physics and the nonlinearity constraint. We then evaluate the training paradigms necessary to optimize these systems, emphasizing physics-aware algorithms that mitigate the simulation-to-reality gap. The discussion proceeds to hardware realization, comparing static and reconfigurable phase-modulation platforms and their respective trade-offs in resolution, bandwidth, and cost. We further isolate wavelength multiplexing as a critical degree of freedom, distinguishing its role as both a chromatic constraint and an enabler of multi-functional computing. Subsequently, we assess the maturity of ODNN applications in imaging and pattern recognition, critically distinguishing between proof-of-concept demonstrations and deployable solutions. The review culminates in a critical examination of unresolved controversies regarding claimed computational advantages and concludes with a roadmap for hybrid electro-optical co-design. By synthesizing these dimensions, this work aims to clarify the current state of optical AI, arguing that while ODNNs offer unique physical advantages, their future success depends on strategic integration with electronic systems to overcome inherent physical limitations.

\section{Physical Architecture of Diffractive Neural Layers: From Diffractive Lenses to Trainable Optical Networks}\label{physical-architecture-of-diffractive-neural-layers-from-diffractive-lenses-to-trainable-optical-networks}

By substituting electronic weight matrices with spatially varying phase profiles that modulate light via Fresnel or Fraunhofer diffraction, Diffractive Deep Neural Networks (D2NNs) fundamentally redefine computation. \citep{doi_10_1126_science_aat8084_e055772d} \citep{doi_10_1117_12_2675468_10b5d458} In this physical paradigm, linear transformations are executed through discrete arithmetic and emerge from wave propagation governed by Fourier optics, where the optical field evolves through free space and phase-modulating layers according to diffraction integrals. Specifically, the complex field at each layer is computed as the Fourier transform of the previous layer's field, modulated by a phase profile, and propagated over an axial distance, effectively approximating matrix-vector multiplication without active power consumption during the forward pass. \citep{doi_10_1117_1_apn_5_1_016005_dfabad74} This architecture addresses the energy and scaling limitations of Moore's Law by leveraging the inherent parallelism of light, enabling real-time, energy-efficient inference for tasks such as biomedical cell classification. While conventional digital systems rely on sequential processing, D2NNs utilize cascaded phase-only modulation planes to perform all-optical signal processing, though their accuracy in simulation may currently trail behind electronic counterparts despite offering distinct advantages in speed and efficiency. Thus, ODNNs represent a distinct physical computer where the laws of diffraction dictate computational expressivity. The reliance on passive diffraction inherently restricts operations to the linear domain, creating a fundamental bottleneck that necessitates the hybrid electro-optical solutions explored later. This limitation sets the stage for understanding how these abstract phase profiles are materially instantiated. \citep{doi_10_1364_ol_593231_ab93febc} \citep{s2_f7b7631268f17b1f868015df0bd08ef57904024d_011f3a59} \citep{doi_10_1063_5_0323145_de1041b1} \citep{doi_10_1109_twc_2025_3632886_01858e99} \citep{doi_10_3390_life13051148_9434b5b6} \citep{doi_10_1364_ol_428761_f1f67ac8} \citep{doi_10_1002_lpor_202200348_d54cae18}

Physical instantiation of Diffractive Neural Network optical weights as spatially varying phase profiles is required to transition from abstract diffraction integral mathematics to functional hardware. This materialization typically relies on multilayer thin-film structures, a photonic platform composed of stacked material layers with thicknesses ranging from tens to hundreds of nanometers. \citep{doi_10_1039_d5nh00367a_1be9a889} \citep{doi_10_1002_smtd_202400087_8e52a90c} In these systems, the specific combination of materials and layer depths is optimized to control light interference, thereby generating the precise phase shifts required for computation. While forward optical responses are straightforward to calculate using methods like the Transfer Matrix Method, determining the optimal physical structure for a desired computational function constitutes a complex inverse design problem. Modern approaches address this by employing deep learning models to map target optical responses directly to structural parameters, effectively automating the creation of these trainable optical substrates. Complementing thin films, all-dielectric metasurfaces offer an alternative architecture consisting of subwavelength-scale optical antennas that manipulate scattered waves through electromagnetic resonances. \citep{doi_10_1088_1612_202x_ab2121_c375efca} \citep{doi_10_1109_mac64480_2025_11140514_7f08c073} These metasurface-based Diffractive Deep Neural Networks (D2NNs) process information by modulating input coherent light fields through a succession of passive diffractive layers, where each layer acts as a collection of individual phase and amplitude modulation units. By inversely optimizing the distribution of these units via error back-propagation, such networks can approximate arbitrary complex-valued linear transformations between input and output fields. However, conventional implementations relying on thickness-based modulation units often face resolution limits comparable to the operating wavelength, highlighting the balance between fabrication simplicity and the miniaturization potential offered by advanced metasurface geometries. \citep{doi_10_1016_j_isci_2025_112222_1f7c10f5} \citep{doi_10_1515_nanoph_2023_0760_2c3ec088}

Single-plane diffractive optical elements (SPDOEs) enable compact imaging systems that overcome the space requirements of traditional multi-lens setups, yet their computational expressivity is fundamentally constrained by the physics of a single diffraction event. In two-dimensional networks with a hard circular pupil, the point spread function is described by a Bessel function of the first kind, and its Fourier transform remains flat over the pupil aperture, establishing a rigid spatial frequency response that limits the complexity of transformations a single layer can perform. To transcend these boundaries, Diffractive Deep Neural Networks (ODNNs) utilize diffraction across cascaded diffractive surfaces to perform arbitrary functions, distinguishing them from their single-layer counterparts. \citep{doi_10_1186_s43593_025_00113_w_ca9e8e1e} These multi-layer architectures provide significantly better performance for generalization and application-specific design, leveraging passive photon propagation to execute machine learning tasks at the speed of light with low energy consumption. However, this enhanced capability introduces practical challenges related to fabrication complexity and acute sensitivity to opto-mechanical alignment; even minor layer-to-layer misalignments can degrade performance unless specifically addressed during the training phase through strategies like `vaccination' against physical variations. This transition from static single elements to trainable cascades sets the stage for understanding how gradient descent is applied directly to physical parameters to optimize these complex optical pathways. \citep{s2_7caf7c29ac69aa12038f926409fb469f65ebeb2a_e9ce584d} \citep{doi_10_1364_oe_579297_3ffb4db0} \citep{doi_10_1364_oe_446890_26d35f03} \citep{doi_10_1515_nanoph_2020_0291_c4a9f42d} \citep{doi_10_1364_optica_576204_30060082}

To realize the universal approximation capabilities of multi-layer cascades, physical implementation necessitates a mechanism for optimizing the phase profiles of individual layers. This is achieved through the differentiability of diffractive propagation, which allows gradient-based training to adjust physical parameters such as meta-atom dimensions and inter-plane spacing directly from input-output data pairs. \citep{doi_10_1117_12_2661856_e8f727e8} \citep{doi_10_1002_mop_70313_0297d592} In practice, forward propagation models ODNN research transmission using approximations like the locally periodic assumption, where each meta-atom introduces a specific phase modulation to the input electric field before near-to-far field transformations propagate the wave to subsequent layers. While early designs often simplified this process by considering only phase modulation with fixed transmission efficiency, the evolution of ODNN research has highlighted the limitations of passive, fixed architectures that cannot be re-trained for new tasks once fabricated. Although programmable nodes using spatial light modulators have been explored to introduce tunability, such reflection-type approaches often limit the number of layers or introduce electronic delays that break light-speed computing capabilities. Consequently, modern trainable architectures increasingly focus on transmission-type configurations where both phase and amplitude modulation can be optimized during training, thereby enhancing computational expressivity while maintaining the ultra-low latency inherent to all-optical systems. \citep{doi_10_1038_s41598_022_19973_0_e156e6e2} \citep{doi_10_1038_s41928_022_00719_9_a2169947}

Leveraging the trainable designs established previously, advanced architectural variants now exploit additional physical degrees of freedom to enhance computational density and functionality. Polarization multiplexing enables a single diffractive network to simultaneously implement a group of complex-valued arbitrary linear transformations at the same output field of view by encoding different tasks into distinct input and output polarization states. \citep{doi_10_1038_s41377_022_00849_x_e31226af} The optical architecture offers significant advantages in speed, versatility, and compactness compared to employing multiple dedicated subsystems, requiring only the addition of linear polarizers without altering the core diffractive architecture. Complementing this, two-dimensional photonic crystal diffractive optical neural networks (PhC-DONNs) realize on-chip parallel classification schemes through multi-wavelength processing, overcoming the inherent limitations of traditional single-wavelength modes by establishing a joint wavelength-diffraction modulation model. However, purely all-optical systems often face challenges related to complex optical setups and the lack of efficient nonlinear activation, motivating the development of hybrid optical-electronic workflows. In these hybrid configurations, optimized diffractive optics can implement color-sensitive kernels to preserve spectral information typically lost in narrow bandpass filtering, while electronic components handle nonlinearities and adaptive training. Such hybridization simplifies the physical architecture and allows for reconfigurable weights without moving parts, bridging the gap between the ultra-low latency of optical propagation and the flexibility of electronic control, thereby advancing the paradigm of light-based neurocomputing. \citep{doi_10_1038_s41377_022_00849_x_e31226af} \citep{doi_10_1515_nanoph_2025_0168_86a0f7aa} \citep{doi_10_1038_s41598_018_30619_y_5f5bfb09} \citep{doi_10_1364_oe_462370_1269b480}

Hybrid architectures expand functional modalities, yet the physical realization of diffractive neural networks remains bound by hard theoretical limits and environmental fragility. The Abbe-Rayleigh diffraction limit imposes a fundamental constraint on spatial resolution, preventing traditional far-field systems from resolving micro- and nano-scale features without resorting to near-field probes that cannot penetrate biological samples or complex superoscillatory designs that rely on precise interference patterns. Beyond these theoretical boundaries, practical implementation faces significant stability challenges; resonators, though advantageous for introducing necessary nonlinear activation functions, are highly vulnerable to temperature fluctuations and mechanical vibrations that degrade performance. Similarly, diffractive elements and interferometers require exact alignment and suffer from noise introduction, making high-precision fabrication costly and operationally difficult. To bridge the gap between idealized mathematical models and these physical realities, Hardware-in-the-Loop (HIL) methodologies have emerged as a critical design strategy. By embedding the actual optical hardware as a black-box component within the optimization loop, HIL approaches account for unmodeled wavefront errors and system discrepancies, enabling the robust design of complex systems like achromatic extended-depth-of-field imagers. Ultimately, acknowledging these constraints is essential for transitioning ODNNs from sensitive laboratory prototypes to reliable computing paradigms where wave physics dictates both capability and limitation. \citep{doi_10_3390_nano14080697_c372a140} \citep{doi_10_1364_oe_461549_b2d41b3f} \citep{doi_10_1117_1_ap_6_5_056004_3f240d06}

Despite the physical constraints and environmental sensitivities detailed previously, diffractive neural networks (D2NNs) offer a distinct computational paradigm defined by ultra-low latency, massive parallelism, and inherent energy efficiency. Unlike electronic systems constrained by the von Neumann bottleneck of data shuttling between memory and processors, D2NNs leverage light propagation to perform matrix-vector multiplications at the speed of light with minimal power consumption, as diffraction and phase modulation occur naturally without additional energy expenditure. This architecture enables high-throughput processing suitable for real-time edge applications, such as all-optical image segmentation for autonomous driving, where direct access to optical degrees of freedom eliminates the need for digital preprocessing. Furthermore, the co-design of hardware structures and software algorithms allows a single physical architecture to support multifunctional operations; by incorporating learned structured illumination or utilizing dual-wavelength differential signals, these networks can switch between diverse tasks like classification and beam steering or capture multiscale spatial features without physical reconfiguration. While traditional single-wavelength designs face limitations in feature extraction, advanced frameworks demonstrate that synergizing complementary optical responses significantly improves accuracy and robustness against noise. Thus, ODNNs function not merely as accelerators but as versatile physical computers where the laws of diffraction dictate a unique balance of speed, parallelism, and multifunctionality. \citep{doi_10_48550_arxiv_2602_07717_16cb1c7a} \citep{doi_10_1364_prj_576203_ae347fba} \citep{doi_10_1038_s41467_022_28702_0_1b79b43c} \citep{doi_10_1002_lpor_202500433_c50e8d3b} \citep{doi_10_1364_oe_598903_731aaac1} \citep{doi_10_1364_ol_591818_8738c4b1}

Leveraging the ultra-low latency and parallelism established in prior sections, diffractive neural networks demonstrate unique utility in specialized imaging and security tasks where physical wave manipulation replaces digital post-processing. In security applications, physics-informed diffractive networks approximate the inverse transmission matrix of scattering media to enable high-speed, all-optical decryption of speckle ciphertexts. By operating directly on the amplitude and phase of the optical field, these systems eliminate the latency inherent in optoelectronic conversion, though their reconstruction quality remains fundamentally constrained by the system's space-bandwidth product and sensitivity to alignment errors. Similarly, in compressive sensing, machine learning algorithms optimize optical architectures using large labeled datasets to perform task-specific classification with significantly fewer detector elements than conventional imaging. While monolithic prism arrays may introduce blurring compared to digital micromirror devices, both architectures achieve superior classification accuracy in low signal-to-noise conditions due to high radiometric throughput. Furthermore, deep learning advancements are accelerating the translation of holography into biomedical contexts, enabling label-free live cell imaging at high frame rates. Although these methods require carefully registered training data to avoid artifacts, the precise control offered by microscopy instrumentation provides an ideal framework for regularizing network convergence. Collectively, these applications underscore that ODNNs are accelerators and distinct physical computers where diffraction laws dictate computational expressivity, trading fixed weights and wavelength dependence for unprecedented speed and energy efficiency in domain-specific tasks. Realizing this potential in practice, however, requires rigorous training methodologies and hardware implementations capable of managing these physical constraints. \citep{doi_10_1117_1_apn_5_1_016010_3e30c230} \citep{doi_10_1117_1_oe_59_5_051404_363a991a} \citep{doi_10_1038_s41377_019_0196_0_c541a42e}

\section{Training Paradigms for Optical Diffractive Networks: Gradient-Based Learning, End-to-End Optimization, and Physics-Aware Algorithms}\label{training-paradigms-for-optical-diffractive-networks-gradient-based-learning-end-to-end-optimization-and-physics-aware-algorithms}

Following the discussion of optical diffractive network applications in the preceding section, it is essential to revisit the foundational training paradigms that enable their physical realization. End-to-end differentiable training establishes the foundational mechanism for optimizing optical diffractive neural networks by treating physical light propagation as a computational graph amenable to gradient descent. \citep{doi_10_1145_3508352_3549378_13c53996} \citep{doi_10_1364_oe_533298_571e0aea} In this framework, the phase profiles of metasurface layers serve as trainable parameters, updated via backpropagation through the scalar diffraction integral that governs forward wave behavior. \citep{doi_10_1117_12_3056838_7df913be} Rather than relying on heuristic design, the system utilizes algorithms like stochastic gradient descent to minimize task-specific loss functions, effectively allowing the network to learn how to modulate light fields directly. That implementation leverages efficient Fourier optics methods, such as the angular spectrum or Rayleigh-Sommerfeld theories, to model interference and diffraction between layers, balancing computational speed with the physical reality of wave propagation. The viability of this framework was decisively demonstrated by the first all-optical diffractive deep neural network, which employed these principles to classify handwritten digits at near-light speeds using 3D-printed phase masks optimized entirely through simulation. By embedding the physics of diffraction directly into the optimization loop, this method bridges the gap between abstract machine learning and tangible optical hardware, setting the stage for evaluating the resulting computational efficiencies. \citep{s2_f7b7631268f17b1f868015df0bd08ef57904024d_011f3a59} \citep{doi_10_3390_nano14221812_4d7d91b1} \citep{doi_10_1038_s41598_022_19973_0_e156e6e2} \citep{doi_10_3390_life13051148_9434b5b6} \citep{doi_10_1364_josab_497148_b087424b}

The Diffractive Deep Neural Network (D2NN) framework builds on the differentiable training foundations established previously by leveraging these optimized physical structures to enable all-optical machine learning, where task-specific information processing occurs entirely through the modulation of coherent light fields across successive passive layers without electronic intermediate steps. This framework exploits the inherent parallelism of light propagation, offering significant advantages in throughput and energy efficiency compared to traditional electronic counterparts by performing computations at the speed of light through free-space diffraction. The differentiable nature of the underlying physics allows for the optimization of millions of artificial neurons via backpropagation, facilitating complex operations such as multi-task learning where partial physical layers are shared to reduce hardware complexity while maintaining robustness. Recent advancements, including dual-wavelength differential networks, demonstrate how leveraging complementary optical responses can significantly improve classification accuracy and robustness against noise, validating the potential for orders of magnitude improvements in computation speed. However, while these systems promise ultra-high density computing for applications ranging from object recognition to logical computing, their practical deployment must eventually address the vulnerability to physical adversarial attacks and the challenges of reconfigurability inherent in static diffractive elements. \citep{doi_10_1515_nanoph_2023_0760_2c3ec088} \citep{s2_88f87b34acdb11cf20ad16d6fbdd5399e8617f8f_a548a720} \citep{doi_10_1002_lpor_202500433_c50e8d3b} \citep{doi_10_1109_dac18074_2021_9586204_0c48190f}

Despite the theoretical efficiency of all-optical inference, a significant simulation-to-reality gap often undermines on-chip schemes, particularly those based on metasurfaces where fabrication errors deviate physical performance from idealized numerical models. To bridge this discrepancy, researchers employ HIL training, a methodology that integrates the physical optical setup directly into the optimization loop rather than relying solely on mathematical propagation models. In the same strategy, the hardware acts as a black-box forward model; for instance, end-to-end designs using programmable phase Spatial Light Modulators optimize hyperparameters by fitting reconstructed images against true targets, effectively absorbing wavefront propagation errors that are difficult to model numerically. Hardware-in-the-loop free space optical computing and optimization of polychromatic multilevel diffractive optics for imaging demonstrates this design by optimizing diffractive elements through physical feedback rather than pure simulation. \citep{doi_10_1109_icton67126_2025_11125227_383c0a4a} However, the success of HIL training depends critically on precise mechanical alignment and energy distribution control. Evaluations of tunable liquid crystal retarders demonstrate that even minor misalignments between lens functions cause noticeable reductions in spot energy and focus, necessitating real-time monitoring of detector patterns to correct low-order aberrations during the optimization process. By embedding these physical constraints directly into the learning cycle, HIL methods ensure that the resulting diffractive networks remain robust against the inevitable imperfections of physical realization. \citep{doi_10_1038_s41598_025_19638_8_1b2db1c9} \citep{doi_10_1364_oe_461549_b2d41b3f} \citep{doi_10_1021_acsphotonics_4c01284_s001_10346917}

Physics-informed optimization strategies explicitly embed fabrication constraints and spectral dispersion directly into the loss landscape to ensure realizability, addressing the need to correct physical errors via hardware-in-the-loop methods. The proposed network treats the inverse design of multilayer thin-film structures not merely as a data-fitting problem, but as a combinatorial optimization where the goal is to identify the best combination of materials and thicknesses that satisfy desired optical responses while adhering to physical limits. By incorporating dispersion formulas, such as Sellmeier equations, and weighting multi-wavelength objectives, these strategies can correct aberrations like Zernike terms while managing the non-trivial complexity of stacking layers with nanometer-scale precision. However, the success of such designs is critically bounded by fabrication variability, which introduces significant uncertainty in broadband diffractive neural network performance. Evidence indicates that at shorter wavelengths, small differences in input wavelengths amplify structural errors, leading to poor classification accuracy when distinguishing closely spaced spectral bands. Furthermore, material properties impose hard physical boundaries; for instance, the solid outer shell of sol-gel polymers requires oil immersion configurations that limit the number of fabricable layers, preventing the depth increases that might otherwise mitigate these errors. To overcome these rigid constraints, data-driven inverse design frameworks utilize differentiable neural simulators to map metasurface geometries directly to their spectral responses. This mechanism replaces traditional, computationally expensive electromagnetic solvers with efficient approximations, enabling the joint end-to-end optimization of both the optical encoder and the computational reconstruction network. \citep{s2_76a83ea3baed75ad88b7f0782eb9a694a330c853_cee16cac} \citep{doi_10_1002_admt_202401999_652aff4b} Such co-optimization has been shown to significantly enhance spectral fidelity and robustness to measurement noise, demonstrating that bridging the simulation-to-reality gap requires not just better training algorithms, but loss functions that inherently respect the physics of fabrication and diffraction. \citep{doi_10_1016_j_isci_2025_112222_1f7c10f5} \citep{doi_10_1002_adpr_202300310_6647aa64} \citep{s2_76a83ea3baed75ad88b7f0782eb9a694a330c853_cee16cac}

Alternative training paradigms address the fundamental breakage of gradient chains caused by discrete device mappings, complementing physics-aware regularization that embeds fabrication constraints into the loss function. Directly inserting these discrete mappings into differentiable optical neural network backpropagation disrupts the optimization process because analog optical devices exhibit non-uniform responses and periodic phase constraints that digital emulations often overlook. To maintain trainability, physics-aware differentiable codesign frameworks utilize techniques like Gumbel-Softmax to enable fully differentiable discrete mappings regardless of the device's number representation format or range. Beyond simulation-based corrections, in situ optical backpropagation implements the learning algorithm directly within the physical system, calculating gradients by measuring forward and backward propagated optical fields based on light reciprocity. The configuration has been demonstrated for diffractive optical neural networks, allowing the network to self-adapt to system imperfections and perform inference at the speed of light, though it requires precise optical measurement capabilities. \citep{doi_10_1364_prj_389553_d7a30418} For scenarios where all-optical complexity is prohibitive, hybrid strategies offer a middle ground; a novel nonlinear generalization framework improves classification accuracy in hybrid networks by integrating electronic layers and using knowledge distillation without requiring additional nonlinear physical elements. Similarly, the Only-train-electrical-to-optical-conversion (OTEOC) strategy simplifies diffractive neural network architectures to a single optical layer with optical readout, intrinsically integrating nonlinear functions during the electronic-to-optical encoding process to avoid slow optical-to-electrical conversions. \citep{doi_10_1364_oe_462370_1269b480} In the context of spectral management for hybrid optical-electronic convolutional neural networks, specific trade-offs arise: while a monochromatic image can be achieved through a narrow bandpass filter, important color information would be lost. To mitigate this, color-sensitive kernels could be optically realized, or chromatic aberrations could be harnessed to extract color information, potentially enhanced by hyperspectral imaging techniques. These diverse methods collectively bridge the gap between theoretical differentiability and physical realizability, ensuring that optical diffractive networks remain computationally effective despite hardware non-idealities. \citep{doi_10_1145_3508352_3549378_13c53996} \citep{doi_10_1364_prj_389553_d7a30418} \citep{doi_10_1364_oe_446890_26d35f03} \citep{doi_10_1364_oe_462370_1269b480} \citep{doi_10_1038_s41598_018_30619_y_5f5bfb09}

Advanced optimization paradigms now exploit specific physical degrees of freedom to create reconfigurable and multi-functional diffractive networks, moving beyond the resolution of fundamental differentiability breaks and non-idealities in training. A primary method is polarization multiplexing, which enables a single static diffractive network to simultaneously perform multiple complex-valued arbitrary linear transformations at the same output field of view by encoding distinct tasks into different input and output polarization states. The resulting design offers significant advantages in speed, versatility, and compactness compared to deploying multiple dedicated subsystems, requiring only the addition of linear polarizers as pre-determined hyperparameters without altering the core diffractive architecture. Beyond polarization, the multi-dimensional light modulation capabilities of metasurfaces facilitate wavelength-multiplexed computing, where recomposable layered structures support parallel operations such as optical encryption and multi-task classification across different spectral channels. While early implementations faced trade-offs between device bulkiness and parallel accuracy, rigorous joint optimization of height maps and pluggable components is unlocking compact, highly parallel processors. Complementing these free-space approaches, two-dimensional photonic crystal diffractive optical neural networks (PhC-DONN) integrate on-chip parallel classification schemes capable of multi-wavelength parallel processing by establishing a wavelength-diffraction joint modulation model. This specific architecture overcomes the inherent limitations of traditional on-chip diffractive networks operating in single-wavelength mode and holds potential for applications in real-time multi-object tracking for autonomous driving, utilizing wavelength-encoded spatial coordinates, as well as three-dimensional reconstruction of pathological slices by employing wavelength dimensions to map depth information. Furthermore, the integration of programmable digital-coding metasurfaces addresses the inherent rigidity of passive fabricated devices, allowing for dynamic control of meta-atom states to implement reconfigurable functions like logic gates and permutation operations in real time. These multiplexing and programmability techniques collectively bridge the gap between fixed simulation models and adaptable physical reality, ensuring that optical neural networks can evolve from single-task demonstrators into versatile platforms capable of handling the high-dimensional complexity required for real-world deployment. \citep{doi_10_1038_s41377_022_00849_x_e31226af} \citep{doi_10_1515_nanoph_2024_0483_f81cb7ab} \citep{doi_10_1038_s41928_022_00719_9_a2169947} \citep{doi_10_1515_nanoph_2025_0168_86a0f7aa}

Leveraging the multiplexing capabilities of reconfigurable architectures, end-to-end differentiable training now scales to high-dimensional tasks that demand robustness against physical disorder and complex scene interpretation. A primary demonstration of this scalability is all-optical inference through unknown random diffusers, where a broadband single-pixel diffractive network learns phase profiles to compensate for scattering without external computing power. The optical architecture achieved a blind testing accuracy of approximately 88.53\% for classifying handwritten digits through entirely new diffusers never encountered during training, significantly outperforming monochrome designs by leveraging wavelength-multiplexed information capacity. Beyond classification, these paradigms extend to pixel-wise segmentation for autonomous driving, utilizing architectures with optical skip connections to process RGB channels directly in the optical domain for accurate lane detection. The versatility of diffractive networks further reaches specialized sensing domains, such as direction-of-arrival estimation for wireless signals, where they modulate electromagnetic waves at the speed of light to overcome the latency and hardware complexity of traditional electronic processors. Distinct from these recognition and sensing tasks, physics-informed diffractive neural networks enable high-speed all-optical decryption of ciphertexts by approximating the inverse transmission matrix of speckle patterns. This capability allows for the direct retrieval of secret images from speckle fields in real-time, effectively eliminating the computational latency caused by optoelectronic conversion and post-processing inherent in conventional approaches. Furthermore, fusing physics priors with deep learning enables compact single-plane imaging systems that reconstruct high-quality images from frequency-spatial features, while polarization-dependent metasurfaces introduce location-awareness to simultaneously distinguish patterns and their spatial positions. These advancements confirm that successful ODNN training relies not merely on backpropagation, but on physics-aware strategies that bridge the simulation-to-reality gap to solve complex, real-world problems. \citep{doi_10_1038_s41377_023_01116_3_aba67a18} \citep{doi_10_48550_arxiv_2602_07717_16cb1c7a} \citep{doi_10_1038_s41377_024_01511_4_2538a27d} \citep{doi_10_1364_oe_579297_3ffb4db0} \citep{doi_10_1117_12_3086084_f1b734fa} \citep{doi_10_1117_1_apn_5_1_016010_3e30c230}

\section{Hardware Realization of Diffractive Neural Networks: Phase Modulation Platforms, Fabrication Techniques, and Reconfigurability Trade-offs}\label{hardware-realization-of-diffractive-neural-networks-phase-modulation-platforms-fabrication-techniques-and-reconfigurability-trade-offs}

The evolution of diffractive neural networks (D2NNs) from theoretical constructs to tangible implementations capitalizes on the inherent parallelism and energy efficiency of light diffraction to execute inference at near-light speeds. Unlike electronic counterparts that suffer from memory bottlenecks during matrix multiplication, these all-optical systems utilize free-space interconnections where computation occurs naturally as light propagates through structured layers. Early demonstrations established this framework by employing terahertz sources and 3D-printed diffractive gratings to successfully classify images, proving that optical computing can achieve ultra-low latency without active energy expenditure during the inference phase. For instance, neural network-based designs have been utilized to create two-focal-spot terahertz diffractive optical elements, demonstrating the capability of such passive systems in specific signal processing tasks. \citep{doi_10_1109_tthz_2023_3344020_a0a550ba} However, the physical realization of D2NNs spans a diverse spectrum of hardware platforms, ranging from static printed phase plates and lithographically fabricated multilevel structures to dynamically reconfigurable spatial light modulators (SLMs) and metasurfaces. While static elements offer simplicity and energy efficiency, dynamic systems using SLMs or digital micromirror devices enable complex, nonlinear network nodes and adaptive learning, albeit with increased system complexity. This fundamental tension between the efficiency of static fabrication and the flexibility of reconfigurable components defines the current hardware landscape, directly influencing whether an ODNN can scale from simulation to robust deployment. \citep{doi_10_1155_2021_6667495_60f04b14} \citep{doi_10_1364_josab_497148_b087424b}

Additive manufacturing has emerged as a primary route for physically realizing diffractive deep neural networks (D2NNs), where 3D-printed phase structures demonstrate experimental performance that aligns closely with numerical simulations. This viability for rapid prototyping is complemented by solution-based chemical methods, such as sol-gel processes, which offer scalable and low-cost alternatives for fabricating large-area porous films with tunable refractive indices, contrasting sharply with the precision but higher complexity of lithographic techniques. The geometric details of coating processes also matter substantially: in directional deposition, the geometry of the substrate induces large differences in the angle of incidence, which directly influences film thickness uniformity and density. Surfaces normal to the incident deposition flux typically exhibit the highest compressive stress levels due to greater thickness and bulk density, whereas off-angle features show lower density and thickness, resulting in reduced interfacial stresses; this variability complicates the design of stress-compensated precision optics and impacts the phase profile fidelity achievable in practical diffractive optical elements. To further reduce manufacturing barriers, minimalist optical diffractive neural networks employ 2-level quantized pixel-wise phase modulation, trading some classification accuracy for significantly reduced fabrication complexity while maintaining robust computational capabilities for tasks like image classification. Although these static methods enable cost-effective scaling to hundreds of millions of neurons using wide-field optical components, they inherently lack the dynamic adaptability of active systems, setting the stage for high-resolution dielectric metasurfaces that address the resolution limits of traditional holography and additive approaches. \citep{doi_10_1126_science_aat8084_e055772d} \citep{doi_10_1038_s41598_025_33607_1_2162f081} \citep{doi_10_1002_lpor_202402303_1f555bd4} \citep{doi_10_1116_1_5011790_9ce3fdef}

Additive manufacturing enables scalable static fabrication but often yields neuron sizes comparable to the wavelength, limiting resolution and introducing high-order diffraction. To overcome these constraints, metasurface-based diffractive neural networks utilize subwavelength pixel resolution to engineer spatial distributions of amplitude and phase with precision unattainable by conventional thickness-based modulators. All-dielectric metasurfaces serve as the physical hardware platform for realizing these networks, processing information via coherent light field modulation through arrays of optical antennas that exploit electromagnetic resonances. This architecture offers an alternative to electronic systems by providing inherent parallelism, energy efficiency, and low latency without requiring external power for phase modulation. By reducing individual neuron sizes to scales significantly smaller than the wavelength, all-dielectric metasurfaces enable fully static implementations that eliminate high-order diffraction and maximize spatial utilization. Representative applications include compact optical logic gates, such as 1-bit half-adders and full-adders, which demonstrate the feasibility of miniaturized, high-density all-optical computing for space-constrained fields like robotics and smart sensors. Specific implementations include all-dielectric metasurface based diffractive neural networks designed to function as 1-bit adders, showcasing their utility in performing basic arithmetic operations optically. \citep{doi_10_1515_nanoph_2023_0760_2c3ec088} \citep{doi_10_3390_nano14221812_4d7d91b1} \citep{doi_10_1515_nanoph_2023_0760_2c3ec088} \citep{doi_10_1117_12_2688132_926a2697}

Reconfigurable systems, such as Spatial Light Modulators (SLMs) and integrated rewritable metasurfaces, enable the dynamic task switching essential for flexible computing, even though static metasurfaces offer efficiency. \citep{doi_10_1002_adom_202300561_c4d7a5b0} \citep{doi_10_32446_0368_1025it_2023_12_35_39_fb1bb2c3} An SLM is a dynamic device used for reconfigurable phase modulation, often serving as a programmable platform for diffractive optical elements. However, deploying these networks reveals a simulation-to-reality gap, where fabrication imperfections and device non-idealities degrade trained performance compared to idealized models. For instance, on-chip rewritable metasurfaces utilizing phase-change materials exhibit errors caused by the fabrication process, such as increased losses or limited spatial control over crystalline rods, which necessitate compensation strategies. To bridge this divide, HIL methodologies incorporate physical measurement feedback directly into the training process. That implementation treats the hardware as a black-box system, allowing end-to-end optimization of parameters like SLM phase patterns without requiring explicit mathematical gradients for wavefront propagation. By using evolutionary strategies to fit reconstructed images to true targets, HIL setups effectively model discrepancies that are difficult to capture numerically. Yet, this adaptability comes with trade-offs; in-situ training can correct errors by erasing and writing diffractive units, but the reliance on iterative physical feedback and the inherent losses of reconfigurable materials highlight the tension between flexibility and operational efficiency that defines the hardware selection landscape. \citep{doi_10_1038_s41598_025_19638_8_1b2db1c9} \citep{doi_10_1364_oe_461549_b2d41b3f}

Bridging the simulation-to-reality gap through physics-aware codesign is critical for the successful deployment of reconfigurable systems, which offer adaptability. This methodology is essential because directly mapping discrete physical device responses into diffractive optical neural networks breaks the gradient chain required for standard backpropagation, as analog optical components exhibit non-uniformities and periodic phase constraints that digital models often overlook. To resolve this, differentiable codesign frameworks employ techniques that enable fully differentiable discrete mappings while respecting the physical limits of hardware, such as the \(2\pi\) periodicity of phase modulation. Beyond discretization, manufacturing imperfections like surface roughness further degrade performance when networks are trained solely on ideal numerical models; physics-aware roughness optimization addresses this by incorporating fabrication errors directly into the training loop to mitigate accuracy loss. Taking this integration further, in situ optical backpropagation allows networks to be trained directly within the optical setup by measuring forward and backward propagated fields based on light reciprocity. The same strategy bypasses the extensive computational overhead of digital algorithms and enables the system to self-adapt to inherent hardware non-idealities, ensuring that the transition from simulation to physical realization maintains both energy efficiency and inference accuracy. \citep{doi_10_1145_3508352_3549378_13c53996} \citep{doi_10_1109_dac56929_2023_10247715_f78d6713} \citep{doi_10_1364_prj_389553_d7a30418}

Beyond fabrication non-idealities, the operational stability of deployed diffractive systems is fundamentally constrained by environmental sensitivity and coherence requirements. Photonic neural network resonators, while offering advantageous nonlinear behavior for implementing activation functions, remain vulnerable to temperature fluctuations and mechanical vibrations that can adversely impact their performance. Similarly, the operational stability of diffractive optical neural networks relies on quasimonochromatic illumination conditions where the coherence length exceeds the maximum path length difference between layers, an assumption that may not hold for real-world applications involving incoherent light. These challenges are compounded in multi-layer architectures; although such designs offer improved generalization, they introduce significant opto-mechanical alignment sensitivity where physical misalignments degrade the computational fidelity of trained phase profiles. To address this, recent strategies employ a ``vaccination'' training scheme that models layer-to-layer misalignments as continuous random variables during optimization, enabling networks to maintain object recognition accuracy despite severe displacement in the x, y, and z directions. Ultimately, selecting a hardware platform requires balancing these stability constraints against application needs, as even robustly trained models must operate within strict physical bounds to transition successfully from simulation to deployment. \citep{doi_10_3390_nano14080697_c372a140} \citep{doi_10_1364_oe_523619_ddeb1b53} \citep{doi_10_1515_nanoph_2020_0291_c4a9f42d}

To meet the demand for operational stability, current advanced integration strategies employ on-chip logic and multiplexing to maximize computational density while maintaining robustness. Metasurface-enabled diffractive neural networks in the visible spectrum utilize multiplexed phase profiles to achieve compact, static hardware realization with high spatial resolution and low power consumption compared to free-space implementations. This miniaturization extends to ultra-compact microscale optical logic operations, where three-dimensional diffractive structures printed via two-photon polymerization can execute multiple logic gates with 100\% accuracy within a volume of just 100 \(\times\) 100 \(\times\) 50 \(\mu\) m\(^{3}\), directly addressing the bandwidth and power limitations of traditional electronic systems. Cascade diffractive architectures further address system constraints by circumventing the need for explicit nonlinear electronic activation functions, relying instead on optical modulation and free-space propagation to achieve computational depth through equivalences between digital and diffractive phase dropout mechanisms. To enhance throughput without mechanical reconfiguration, multi-wavelength architectures encode distinct tasks into separate spectral channels, allowing simultaneous classification of different databases while preserving individual task accuracy. For instance, diffractive deep neural networks have been employed for 3D holographic imaging by leveraging wavelength division multiplexing to process distinct spectral channels simultaneously. \citep{doi_10_1109_acp_ipoc63121_2024_10809872_a7935774} Complementing spectral diversity, orbital angular momentum (OAM) mode multiplexing introduces a depth-controllable dimension by combining OAM modes with spatial depths through unitary transformations, enabling the reconstruction of multiple holograms within a single five-layer network. These designs collectively demonstrate that integrating multiplexing dimensions with on-chip fabrication transforms diffractive hardware from single-task prototypes into high-capacity, multitasking processors, setting the stage for selecting the optimal platform based on specific application constraints. \citep{doi_10_3390_mi16010008_4be72643} \citep{doi_10_1021_acsphotonics_5c02472_283a8bea} \citep{doi_10_1515_nanoph_2022_0615_9a3be8d6} \citep{doi_10_1364_oe_447337_61ffab71} \citep{doi_10_1364_ol_428761_f1f67ac8}

Synthesizing the hardware landscape reveals that platform selection is fundamentally a constrained optimization between static efficiency and dynamic adaptability. Conventional metasurfaces function as static, `write-once' devices where optical properties are fixed upon fabrication, a characteristic that severely restricts their flexibility in dynamic scenarios despite offering strong wavefront shaping capabilities. Similarly, current on-chip diffractive neural network architectures become fixed once their physical structures are trained and fabricated, rendering them unable to adapt to new tasks without costly re-fabrication. While these static platforms enable diffractive deep neural networks to achieve orders of magnitude improvements in computation speed and energy efficiency over electronic counterparts through the parallel nature of optical propagation, this performance comes with critical vulnerabilities. Recent analyses demonstrate that pre-trained models can be susceptible to physically feasible adversarial attacks, such as phase or amplitude perturbations, which can drastically reduce classification accuracy. Consequently, while static platforms maximize energy efficiency for dedicated tasks, the lack of post-fabrication reconfigurability remains a significant bottleneck for applications requiring robustness against environmental changes or adversarial threats. \citep{doi_10_1002_nap2_70185_56c8ab30} \citep{doi_10_1117_12_3109287_f4650a31} \citep{doi_10_1109_dac18074_2021_9586204_0c48190f}

\section{Wavelength Multiplexing and Spectral Bandwidth in Diffractive Optical Networks: Multi-Wavelength Computation and Chromatic Constraints}\label{wavelength-multiplexing-and-spectral-bandwidth-in-diffractive-optical-networks-multi-wavelength-computation-and-chromatic-constraints}

Optical neural networks based on diffraction (ODNNs) operate as inherently dispersive systems in which the computational transfer function is strictly coupled to the illumination wavelength. This chromatic dispersion arises because the phase modulation imparted by a fixed physical structure scales inversely with wavelength, causing the point spread function and diffraction angles to shift as the input spectrum deviates from the design condition. \citep{doi_10_1117_1_oe_65_5_058106_f60e3d89} Consequently, a network optimized for a specific monochromatic source experiences a drift in its linear transformation matrix when operated under broadband conditions, leading to degraded computational accuracy unless narrow bandpass filtering is employed to isolate color channels. While this physical dependency imposes strict bandwidth limits on single-task operation, it simultaneously establishes the foundational mechanism for wavelength multiplexing, where distinct spectral components can trigger independent computational pathways within the same static volume. Understanding this dual nature of dispersion is critical, as it defines the boundary between performance loss in uncontrolled environments and the opportunity to harness spectral degrees of freedom for multi-functional optical computing. \citep{doi_10_1038_s41598_018_30619_y_5f5bfb09} \citep{doi_10_3390_nano14221812_4d7d91b1} \citep{doi_10_1109_icton67126_2025_11125227_383c0a4a}

The physics of chromatic dispersion establishes that wavelength shifts alter phase delays, yet the practical consequence for diffractive neural networks trained at a single wavelength is a severe degradation in inference capability when exposed to spectral deviations. Empirical studies reveal that attempting to classify multiple wavelengths in parallel using shallow architectures often results in poor performance; for instance, single- and two-layer networks designed to distinguish three visible wavelengths achieved accuracies of only 53\% and 60\%, respectively, frequently misclassifying distinct inputs as a dominant intermediate wavelength. This failure stems not only from the inherent coupling of modulation functions across the spectrum but also from fabrication constraints that limit the number of diffractive layers available to resolve complex spectral features. In sharp contrast, when scattering conditions remain strictly consistent with the monochromatic training environment, these same single-wavelength optimized systems can maintain functional utility, achieving blind testing accuracies exceeding 78\% even when imaging through unknown random diffusers. However, this robustness is fragile; comparative analyses indicate that such monochrome designs perform significantly worse than their broadband counterparts under identical challenging conditions, underscoring that the high accuracy of single-wavelength networks is contingent upon a narrow spectral match that rarely exists in practical, broadband illumination scenarios. \citep{doi_10_1002_adpr_202300310_6647aa64} \citep{doi_10_1038_s41377_023_01116_3_aba67a18}

Published studies have developed achromatic diffractive networks that employ multi-layer architectures and hardware-aware optimization to stabilize phase responses across broad spectra, aiming to counteract the performance degradation caused by chromatic dispersion. A primary strategy involves integrating multi-layer thin films with dual-scale nanostructures, where a multi-wavelength weighted optimization framework collaboratively minimizes reflectance and corrects wavefront aberrations; this design has successfully extended spectral bandwidths to 140 nm while suppressing low-order errors to mere milliwavelengths. Alternatively, bilayer metalens architectures and 3D-printed multilayer structures offer a geometric route to expand operational bandwidth beyond the limits of single-layer counterparts. When fabrication variability threatens performance---particularly at shorter wavelengths where small dimensional errors disproportionately affect inference accuracy---hybrid designs employing hardware-in-the-loop (HIL) optimization provide a robust solution. By embedding a programmable spatial light modulator within the training loop, HIL methods bypass inaccurate mathematical models of wavefront propagation, directly optimizing phase patterns against physical reality to achieve achromatic extended depth-of-field imaging. In the context of hybrid diffractive optics, hardware-in-the-loop methodologies have been employed to design systems for achromatic extended-depth-of-field imaging. By integrating a programmable spatial light modulator directly into the training loop, the proposed network allows phase patterns to be optimized against physical reality rather than relying on potentially inaccurate mathematical models of wavefront propagation. \citep{doi_10_1364_oe_461549_b2d41b3f} However, these mitigation strategies face tangible boundaries: increasing the number of diffractive layers to improve multi-wavelength classification is often constrained by polymer properties and the working distance of high-numerical-aperture objectives, which limit the total fabricable height and resolution. Thus, while engineering solutions can significantly suppress dispersion, they must navigate strict trade-offs between structural complexity and manufacturability to maintain the high-efficiency broadband operation required for practical optical computing. \citep{doi_10_1002_adpr_202300310_6647aa64} \citep{doi_10_3390_modelling7010029_9851fe61} \citep{doi_10_1364_oe_461549_b2d41b3f} \citep{doi_10_29026_oea_2021_200008_8948c2c3} \citep{doi_10_1126_sciadv_adj9262_15884f2e}

Rather than suppressing chromatic dispersion, diffractive neural networks can exploit optical degrees of freedom to implement multiple independent computational functions within a single physical structure. This strategy, known as wavelength multiplexing, enables a static diffractive element to perform distinct tasks simultaneously by encoding separate functions for different spectral bands. The mechanism relies on optimizing the phase profile of cascaded layers so that specific wavelengths trigger unique linear transformations, effectively mapping input degrees of freedom to targeted outputs without mechanical reconfiguration. While traditional designs often partition devices into separate regions that limit efficiency when illuminated by a single wavelength, multi-layer adaptive designs overcome this by training the entire aperture to focus or classify correctly at multiple target wavelengths concurrently. For instance, research demonstrates that assigning different wavelengths as basebands or carrier frequencies allows a network to balance learning rates and execute multitask operations, such as parallel image recognition, with higher efficacy than single-task baselines. Wavelength division multiplexing has been utilized in 3D holographic imaging with diffractive deep neural networks to leverage chromatic dependence as a functional resource. This method enables distinct wavelengths to trigger independent operations within the same physical volume, facilitating complex tasks such as 3D holographic reconstruction through wavelength-specific signal processing. \citep{doi_10_1109_acp_ipoc63121_2024_10809872_a7935774} The configuration transforms the inherent spectral sensitivity of diffractive systems from a liability into a resource for high-throughput computing, though it requires careful hardware-software co-design to manage task accuracy trade-offs. Ultimately, leveraging wavelength dependence in this manner aligns with the chapter's thesis that dispersion provides a critical degree of freedom for parallel, multi-functional optical computing. \citep{doi_10_1038_s41377_022_00903_8_098356a7} \citep{doi_10_1364_prj_576203_ae347fba} \citep{doi_10_1364_ol_460186_da35c128} \citep{doi_10_1155_2020_9748380_13fcb0d2} \citep{s2_88f87b34acdb11cf20ad16d6fbdd5399e8617f8f_a548a720}

Polarization multiplexing further expands the capacity of a single diffractive network to implement multiple complex-valued linear transformations concurrently within the same output field of view, building on the use of wavelength as a computational degree of freedom. The resulting design leverages the vector nature of light, where distinct input and output polarization states act as independent channels for programming tasks such as object detection or classification without altering the underlying diffractive architecture. Evidence demonstrates that integrating these tasks into one system offers significant advantages in speed, compactness, and cost-effectiveness compared to deploying multiple dedicated subsystems, requiring only the addition of readily available linear polarizers as pre-determined hyperparameters between layers. To achieve even higher information density, orbital angular momentum (OAM) mode multiplexing can be combined with wavelength and polarization dimensions, enabling the simultaneous reconstruction of multiple holograms through depth-controllable spatiotemporal evolution. By exploiting multi-plane light conversion within a multi-layer optical diffractive neural network, experimental reports have successfully modulated both OAM modes and spatial depths to realize deep multiplexing holography with high fidelity. Researchers demonstrated orbital angular momentum deep multiplexing holography by employing an optical diffractive neural network to modulate both OAM modes and spatial depths, achieving high-fidelity results. \citep{doi_10_1364_oe_447337_61ffab71} These advanced schemes illustrate that chromatic dependence is a constraint to be mitigated and a foundational element of a broader strategy where multiple physical dimensions are fused to maximize the functional throughput of static optical processors. \citep{doi_10_1038_s41377_022_00849_x_e31226af} \citep{doi_10_1364_oe_447337_61ffab71}

Realizing the multi-functional potential of these high-density multiplexed systems, where precise control over wavelength and polarization channels is paramount, requires specialized training paradigms to bridge the gap between idealized simulations and physical hardware constraints. Physics-aware differentiable discrete codesign addresses this by maintaining the gradient chain during backpropagation despite the non-uniform, analog nature of real optical devices. Unlike standard quantization, this framework employs techniques such as Gumbel-Softmax to enforce physical boundaries, ensuring that trainable phase parameters respect periodic limits (e.g., restricted within {[}0, 2 \(\pi\) {]}) while accommodating device-specific non-monotonic responses. \citep{doi_10_1145_3508352_3549378_13c53996} Complementing this, neural network-based design methods utilize convolution simulations of the diffraction integral to optimize diffractive elements for specific spectral ranges, such as the terahertz band. These algorithms can generate unintuitive phase maps that redirect radiation into complex intensity distributions, offering solutions beyond traditional deterministic holography. Ultimately, the design of these layers involves iteratively adjusting complex-valued transmission coefficients for statistical tasks, a process that inherently ties the network's final functionality to the illumination spectrum used during training. By embedding these physical realities directly into the optimization loop, researchers can create robust static structures capable of executing precise multi-wavelength computations. \citep{doi_10_1145_3508352_3549378_13c53996} \citep{doi_10_1109_tthz_2023_3344020_a0a550ba} \citep{doi_10_1021_acsphotonics_0c01583_7c90196c}

Physics-aware optimization establishes the structural blueprint for diffractive networks, yet enhancing their computational expressivity and robustness requires addressing inherent linearity constraints and environmental sensitivities. One strategy involves integrating nonlinear photonic components to implement activation functions directly within the optical domain; resonators, for instance, offer advantageous nonlinear behavior for this purpose but introduce vulnerability to temperature fluctuations and mechanical vibrations that can degrade performance. In contrast, interferometers excel at manipulating phase and amplitude for complex computations yet remain susceptible to phase variations requiring precise alignment. To mitigate such instabilities without relying solely on fragile hardware nonlinearities, ensemble learning strategies have been applied to passive diffractive deep neural networks (D\(^{2}\)NNs). By diversifying base models through spatial feature engineering and exploiting parallel processing across multiple independently trained networks, the optical architecture significantly improves statistical inference and generalization capabilities beyond what single deterministic layers achieve. These methods collectively transform the static physical structure from a brittle, single-task processor into a more resilient computing substrate, balancing the balance between enhanced functionality and the practical challenges of fabrication precision and environmental stability before considering broader system-level applications. \citep{doi_10_3390_nano14080697_c372a140} \citep{doi_10_1038_s41377_020_00446_w_c890b40c}

Balancing inherent physical advantages against current accuracy limitations is ultimately required for deploying diffractive optical neural networks in high-stakes environments. While digital convolutional neural networks may still outperform optical systems in raw classification metrics, all-optical diffractive classifiers offer distinct benefits in real-time, energy-efficient inference, making them attractive for applications where speed and power consumption are critical constraints. This trade-off is evident in emerging proposals for autonomous driving, where diffractive architectures aim to improve the latency and computational efficiency of large-scale digital models by enabling direct, all-optical processing of RGB scenes without digital preprocessing. In the context of autonomous driving, all-optical segmentation via diffractive neural networks has been proposed to process RGB scenes directly without digital preprocessing, aiming to enhance latency and computational efficiency. \citep{doi_10_48550_arxiv_2602_07717_16cb1c7a} The hardware motivation for this shift is reinforced by integrated diffractive photonic computing units, which provide ultra-high bandwidth and massive parallelism compared to electronic counterparts, although scaling these systems often faces challenges related to footprint and energy consumption in programmable components. Historically, the first all-optical diffraction deep neural network systems demonstrated these ultra-low latency and low power advantages using passive 3D-printed layers, yet they operated within strict constraints defined by their static hardware implementation. To overcome such rigidity while maintaining efficiency, recent advancements introduce multifunctional diffractive processors with learned structured illumination; by incorporating illumination patterns as trainable parameters, these systems enable the reuse of static diffractive layers to switch rapidly between independent tasks, such as classification and computational imaging. Thus, while broadband dispersion remains a challenge for single-task accuracy, it simultaneously provides the degrees of freedom necessary for multi-functional optical computing, positioning diffractive networks as a transformative technology for label-free biomedical analysis and resource-constrained edge computing. \citep{doi_10_1364_ol_593231_ab93febc} \citep{doi_10_48550_arxiv_2602_07717_16cb1c7a} \citep{doi_10_1038_s41467_022_28702_0_1b79b43c} \citep{doi_10_1364_josab_497148_b087424b} \citep{doi_10_1364_ol_591818_8738c4b1}

\section{Applications of Diffractive Neural Networks: Imaging, Classification, Segmentation, and Pattern Recognition}\label{applications-of-diffractive-neural-networks-imaging-classification-segmentation-and-pattern-recognition}

The all-optical machine learning framework known as Diffractive Deep Neural Networks (D2NNs) performs inference at the speed of light by encoding neural network weights into the physical thickness of diffractive layers. \citep{doi_10_1021_acsphotonics_4c01284_s001_10346917} Unlike electronic processors constrained by data transfer bottlenecks and transistor-induced latency, D2NNs utilize wave propagation through free space to execute analogue computing instantaneously as signals move through the system. This architecture has been experimentally validated for multi-class image classification, specifically demonstrating high accuracy on standard benchmarks like the MNIST handwritten digit dataset. In one prominent implementation operating at visible wavelengths, a five-layer D2NN fabricated on a SiO2 substrate achieved a blind classification accuracy of 91.57\% on held-out test sets without requiring any electronic computation during the inference phase. Beyond numeric digits, this capability extends to alphabetic image recognition, where multi-channel diffractive designs have demonstrated significant improvements in accuracy compared to traditional single-channel architectures when classifying distinct letter datasets. Further expanding this capability, recent studies have developed nonlinear all-optical implementations using specific wavelengths, such as 10.6 \(\mu\) m, alongside visible-light metasurface systems capable of complex object recognition. While these demonstrations establish a foundational maturity for single-function tasks, they primarily serve as proof-of-concept benchmarks; the transition from these controlled optical environments to general-purpose computing remains contingent on overcoming significant simulation-to-reality gaps and environmental sensitivities. \citep{doi_10_54097_hset_v24i_3957_e6f20e8a} \citep{doi_10_1016_j_eng_2020_07_032_6e23cd05} \citep{doi_10_1155_2021_6667495_60f04b14} \citep{doi_10_1109_piers53385_2021_9694760_15b5fd61}

The fundamental physical advantages of diffractive optical neural networks (DONNs) regarding latency and energy consumption serve as the primary motivation for their adoption, extending beyond the classification benchmarks established in standard datasets. Unlike electronic circuits that rely on electron transport and resistive heating, DONNs perform computations passively using photons, enabling inference at the speed of light with significantly reduced power requirements. This photonic hardware acceleration offers ultra-high bandwidth and massive parallelism, addressing the resource bottlenecks inherent in scaling artificial neural networks on conventional silicon platforms. Research indicates that diffractive deep neural networks (D2NNs) have demonstrated orders-of-magnitude improvements in both computation speed and energy efficiency compared to their electronic counterparts, establishing their viability for high-speed, low-power applications. \citep{doi_10_48550_arxiv_2305_03627_fb0d0b27} \citep{s2_e9e6fc7353675c8e70b11e4ba8a931d959c2e912_909a5a0c} However, while integrated photonic circuits provide a promising path toward compact, CMOS-compatible computing units, challenges remain in space utilization and the energy costs associated with controlling phase modulators. These physical efficiencies position ODNNs as compelling solutions for specific high-throughput tasks, yet they also highlight the trade-offs between passive optical speed and the control complexity required for programmable deployment. \citep{doi_10_1002_lpor_202200348_d54cae18} \citep{doi_10_1109_dac18074_2021_9586204_0c48190f} \citep{doi_10_1038_s41467_022_28702_0_1b79b43c}

To expand computational density beyond simple intensity modulation, multiple studies have developed multiplexing strategies that leverage the latency and energy advantages of diffractive systems. Polarization multiplexing is a technique that enables a single diffractive network to perform multiple complex-valued arbitrary linear transformations simultaneously at the same output field of view by encoding data into different input and output polarization states. \citep{doi_10_1038_s41377_022_00849_x_e31226af} That implementation allows various machine vision tasks, such as detection and classification, to be integrated within one compact system, offering greater versatility and cost-effectiveness compared to employing multiple dedicated subsystems. Unlike spatial or wavelength-division methods that may require architectural changes, polarization multiplexing primarily requires the addition of linear polarizers, which can be flexibly coupled with other degrees of freedom to further increase capacity. Complementing this, Orbital Angular Momentum (OAM) mode multiplexing provides a distinct strategy for reconstructing multiple holograms by exploiting the independence of physical dimensions involving wavelength and polarization. By designing optical diffractive neural networks that modulate both OAM modes and spatial depths through unitary transformation, it becomes possible to realize depth-controllable spatiotemporal evolution for high-capacity holography. Furthermore, OAM-encoded diffractive networks have demonstrated the ability to perform all-optical parallel classification on datasets like MNIST using OAM spectrum analysis, though this currently demands interferometer detectors with high sensitivity to manage signal losses. While these multiplexing methods significantly enhance information capacity, they introduce trade-offs between classification accuracy and output efficiency, highlighting that while the theoretical bandwidth is vast, practical deployment still faces challenges in detector robustness and signal integrity. \citep{doi_10_1038_s41377_022_00849_x_e31226af} \citep{doi_10_1364_oe_447337_61ffab71} \citep{doi_10_1117_1_apn_2_6_066006_d88755f0}

Beyond generic classification, diffractive neural networks excel in specialized imaging tasks that address fundamental physical constraints, such as the diffraction limit and chromatic aberration. To achieve superresolution, these systems utilize optical super-oscillations, which are rapid sub-wavelength spatial variations generated by the precise interference of coherent waves to enable far-field imaging beyond the traditional Abbe-Rayleigh limit without the phototoxicity associated with fluorescence microscopy. \citep{doi_10_1117_12_2576293_e7634e1d} \citep{doi_10_1109_cleo_europe_eqec57999_2023_10231830_c702ffd4} Similarly, monolithic diffractive elements tackle achromatic imaging by relaxing trade-offs between aperture size and bandwidth; for instance, meta-axicon clusters generate wavelength-independent Bessel beams to ensure wideband consistency, though this often requires sacrificing focusing efficiency for improved image restoration. A monolithic integrated meta-axicon cluster addresses achromatic imaging constraints by converting oblique plane waves into wideband uniform off-axis Bessel beams. This design, comprising nine meta-axicons with different field angles, enables a meta-camera to achieve achromatic imaging within a stitched 10 \(^\circ\) field of view with angular resolution approaching that of near-diffraction-limit lenses. \citep{doi_10_1038_s41377_026_02272_y_42e58df2} Pushing maturity further, hybrid designs employ hardware-in-the-loop methodologies where programmable spatial light modulators are optimized end-to-end with electronic image signal processing. The same strategy treats the optical hardware as a black box, using zero-order evolutionary strategies to correct for physical discrepancies that pure mathematical models miss, thereby enabling robust extended-depth-of-field imaging. \citep{doi_10_1364_oe_461549_b2d41b3f} These applications demonstrate that while ODNNs offer unique solutions to specific optical bottlenecks, their most effective implementations often rely on co-designing optical preprocessing with computational post-processing. \citep{doi_10_1117_1_ap_6_5_056004_3f240d06} \citep{doi_10_1038_s41377_026_02272_y_42e58df2} \citep{doi_10_1364_oe_461549_b2d41b3f}

Hybrid optical-electronic convolutional neural networks (CNNs) have emerged as a pragmatic architecture that pushes initial computation into passive optics while retaining electronic flexibility for complex decision-making, aiming to bridge the gap between theoretical efficiency and practical deployment. In this design, optimized diffractive elements handle color-sensitive kernels or front-end feature extraction directly in the optical domain, addressing the information loss inherent in monochromatic filtering and relieving electronic hardware of significant computational loads. This division of labor is particularly valuable for latency-critical applications like autonomous vehicles, where optical processing offers speed advantages without the power demands of purely digital models. Complementing this hybrid strategy, compressive sensing systems leverage machine learning algorithms to minimize data collection requirements, utilizing task-specific optical components such as prism arrays or digital micromirror devices to implement sensing matrices. These systems demonstrate that even with optical blurring or lower signal-to-noise ratios, they can achieve classification accuracy comparable to traditional imaging while requiring drastically fewer detector elements. Furthermore, parallel coded ptychography systems illustrate how diffractive principles can overcome scaling obstacles in biomedical microscopy, employing disorder-engineered surfaces to achieve high-throughput, turnkey imaging that eliminates the need for precise mechanical scanning. Distinct from these imaging and classification paradigms, physics-informed diffractive neural networks have been applied to all-optical security schemes capable of real-time decryption of speckle ciphertexts. By approximating the inverse transmission matrix of a scattering medium, these networks enable the direct retrieval of secret images from speckle fields in the optical domain, thereby eliminating the computational latency caused by optoelectronic conversion and post-processing found in conventional approaches. While these architectures successfully mitigate specific limitations regarding power, data volume, and mechanical stability, they often rely on bulk optics or specialized fabrication that highlights the ongoing trade-off between system compactness and performance, setting the stage for the development of reconfigurable materials that can dynamically update these static optical weights. \citep{doi_10_1038_s41598_018_30619_y_5f5bfb09} \citep{doi_10_1117_1_oe_59_5_051404_363a991a} \citep{doi_10_1021_acsphotonics_1c01085_479fdea3} \citep{doi_10_1117_1_apn_5_1_016010_3e30c230}

Emerging hardware platforms utilize reconfigurable materials that allow post-fabrication weight updates to overcome the static limitations of fixed diffractive arrays. Digitized phase-change material heterostacks enable transmissive diffractive optical neural networks with reconfigurability, supporting high-throughput and energy-efficient machine learning processing for applications like computer vision by leveraging the extreme parallelism and low static energy consumption of photons. Alternatively, magneto-optical diffractive deep neural networks utilize Monte Carlo optimization with random domain updates and weight flips to enable real-time, reconfigurable learning in magnetic-optical layers, where candidate configurations are accepted or rejected according to a Metropolis criterion to reduce loss without derivatives. Distinct from these designs, two-dimensional photonic crystal diffractive optical neural networks (PhC-DONNs) realize on-chip optical parallel classification schemes with multi-wavelength parallel processing capabilities. By establishing a wavelength-diffraction joint modulation model, this architecture overcomes the inherent limitations of traditional on-chip diffractive networks operating in single-wavelength mode, demonstrating exceptional performance in image and video classification tasks. Despite these advances, on-chip diffractive optical neural networks based on integrated rewritable metasurfaces face a simulation-to-reality gap that limits their immediate deployment viability compared to electronic baselines. This discrepancy arises from fabrication errors, though the rewritable nature of these systems facilitates in-situ training procedures to compensate for phase and power deviations. While these mechanisms address the rigidity of earlier designs, the persistence of physical imperfections suggests that reconfigurability alone does not fully resolve the challenges required for general-purpose deployment. \citep{doi_10_1002_adpr_202400201_566e9e2f} \citep{doi_10_1051_jeos_2026026_1bd8ff13} \citep{doi_10_1038_s41598_025_19638_8_1b2db1c9} \citep{doi_10_1515_nanoph_2025_0168_86a0f7aa}

Leveraging the potential for dynamic hardware updates, multifunctional diffractive processors now enable single-device multi-task learning, allowing a physical system to execute distinct inference operations without retraining or mechanical reconfiguration. This capability relies on hardware-software co-design where illumination patterns or metasurface arrangements are treated as trainable parameters during optimization. For instance, recent frameworks incorporate ``learned'' structured illumination into the network design, enabling the reuse of fixed diffractive layers while rapidly switching between tasks such as MNIST classification, Fashion-MNIST recognition, and computational imaging. Similarly, metasurface-empowered architectures leverage ultrahigh neuron density and parallelism to support freely arrangeable multi-task functions, effectively addressing the memory-processor bottleneck found in conventional electronic systems. By physically implementing deep learning through light diffraction, these systems demonstrate that a single optical platform can handle diverse applications ranging from visual recognition to beam shaping. However, while these proof-of-concept demonstrations highlight significant flexibility, they also underscore the current transition phase where specialized multi-functionality is achievable, yet general-purpose deployment remains constrained by the very environmental sensitivities and simulation-to-reality gaps that define ODNN research's ongoing challenges. \citep{doi_10_1364_ol_591818_8738c4b1} \citep{doi_10_1364_prj_576203_ae347fba}

Despite the demonstrated capacity for multi-task learning, the transition of diffractive deep neural networks from laboratory prototypes to widely adopted technologies is constrained by a persistent gap between theoretical design and physical robustness. While these systems are celebrated for ultra-low latency and highly parallel optical computing capabilities, their practical deployment faces significant hurdles regarding environmental stability and fabrication precision. Photonic architectures often leverage nonlinear behaviors to implement essential activation functions, yet specific components like resonators remain acutely vulnerable to external perturbations such as temperature fluctuations and mechanical vibrations, which can degrade performance. Furthermore, while diffractive elements offer scalable parallel processing, they introduce noise and aberrations that compromise computational accuracy, and achieving the necessary high-precision fabrication remains both challenging and costly. Although research has validated the viability of all-optical classification for specialized diagnostic tasks, such as identifying COVID-19 from chest X-rays using ensemble deep learning models, these successes currently highlight a niche application space rather than general-purpose readiness. Consequently, while the foundational potential for near-light-speed inference is established, widespread adoption is currently limited by the need for sophisticated control mechanisms to ensure stability against phase variations and the ongoing difficulty of integrating robust, miniaturized photonic circuits that can operate reliably outside controlled environments. \citep{doi_10_3390_nano14080697_c372a140} \citep{doi_10_3390_app122010535_b538921c} \citep{doi_10_1364_josab_497148_b087424b}

\section{Controversies and Unresolved Debates in Diffractive Neural Network Research}\label{controversies-and-unresolved-debates-in-diffractive-neural-network-research}

A central debate in diffractive neural network research questions whether these systems provide a real computational edge or simply serve as sophisticated engineering showcases of established optical principles. Proponents argue that analog optical computing provides inherent parallelism, energy efficiency, and low latency, making it particularly capable of handling massive computational tasks like object classification and machine vision where digital systems face bottlenecks. This perspective is supported by frameworks such as Diffractive Deep Neural Networks (D2NN), which utilize deep learning algorithms to engineer light-matter interactions across passive diffractive surfaces, enabling task-specific inference at the speed of light without continuous electronic power consumption. \citep{doi_10_1021_acsphotonics_4c01284_s001_10346917} Furthermore, integrated chip-based approaches aim to overcome the quadratic scaling of footprint and energy consumption seen in traditional silicon photonic circuits by leveraging ultra-high bandwidth and high-density integration. However, critics contend that these theoretical gains are often offset by the intricate, time-consuming design processes and substantial fabrication complexities required to achieve precise modulation capabilities. ODNN research currently lacks consensus on this trade-off, as reported performance metrics frequently depend on specific benchmark datasets and simulation conditions rather than fundamental physical limits, with some analyses suggesting that the absence of nonlinear activation functions in purely optical layers can hinder optimization when paired with powerful electronic back-ends. \citep{doi_10_3390_nano14221812_4d7d91b1} \citep{doi_10_1515_nanoph_2022_0358_04e3cc8d} \citep{doi_10_1038_s41467_022_28702_0_1b79b43c} \citep{doi_10_1109_jstqe_2019_2921376_c2bef7d6}

Validating claims of computational advantages is hindered by a critical barrier known as the simulation-to-reality gap, which represents the discrepancy between performance predicted by idealized scalar diffraction models and actual hardware results due to fabrication tolerances and material non-idealities. \citep{doi_10_1117_12_3013482_d44587c1} \citep{doi_10_1117_12_3084472_3a4f4f5d} While numerical simulations using methods like the Transfer Matrix Method allow for straightforward prediction of optical responses, the inverse design process often relies on assumptions that do not hold in physical prototypes. Evidence from on-chip implementations indicates that this gap arises specifically from errors introduced during the fabrication process, necessitating algorithmic compensation strategies such as phase and power correction to restore functionality. Furthermore, standard training protocols that assume perfect layer alignment fail to account for the mechanical realities of assembly; however, introducing random misalignment variables during the training phase, a technique termed ``vaccination,'' can significantly improve error tolerance margins against severe layer-to-layer displacements. \citep{doi_10_1515_nanoph_2020_0291_c4a9f42d} The reliance on linear materials in many experimental demonstrations also introduces optical losses that must be carefully managed to maintain signal-to-noise ratios, contrasting with lossless theoretical models. Consequently, a significant portion of the reported performance disparity between diffractive and electronic systems may reflect these favorable simulation assumptions rather than inherent physical limitations, suggesting that claimed energy and speed benefits remain overstated until rigorous experimental validation accounts for scattering, absorption, and manufacturing variances. \citep{doi_10_1016_j_isci_2025_112222_1f7c10f5} \citep{doi_10_1038_s41598_025_19638_8_1b2db1c9} \citep{doi_10_1515_nanoph_2020_0291_c4a9f42d} \citep{doi_10_1126_science_aat8084_e055772d} \citep{s2_94f3822c453a2ab6d44c59f1f60026acb15739af_1ae3e4a4}

End-to-end differentiable training defines a paradigm where backpropagation flows directly through the diffraction integral to optimize optical layers, theoretically enabling high-throughput all-optical processing. However, the proposed network encounters a fundamental gradient chain break when discrete device mappings are introduced, as physical discretization disrupts the smooth differentiability assumed by standard algorithms. While idealized simulations suggest that such networks can scale to hundreds of billions of connections using passive diffractive layers, real-world deployment reveals a mismatch between numerical modeling and physical hardware due to non-uniform device responses and analog constraints. Specifically, optical devices exhibit non-monotonic behaviors and periodic phase limits that uniform digital computation fails to capture, necessitating physics-aware frameworks like Gumbel-Softmax to enable flexible, differentiable discrete mappings. \citep{doi_10_1145_3508352_3549378_13c53996} Furthermore, because physical diffractive layers lack native nonlinearity, reliance on in-silico training with assumed activation functions creates a simulation-reality gap that often requires hybrid electronic assistance or complex fabrication to approximate. These limitations highlight that while end-to-end training offers theoretical speed benefits, its practical superiority remains contingent on resolving the disconnect between idealized gradient chains and the irregularities of analog optical hardware. \citep{doi_10_1145_3508352_3549378_13c53996} \citep{doi_10_1126_science_aat8084_e055772d} \citep{doi_10_1109_dac56929_2023_10247715_f78d6713} \citep{doi_10_1126_sciadv_adu0904_786138ff}

ODNN research now faces a fundamental architectural divide between single-function optimized designs and general-purpose trainable systems. Single-function architectures, such as chromatic-aberration-corrected diffractive lenses, utilize iterative algorithms to optimize groove heights for specific broadband focusing tasks, achieving superior performance for that singular application but lacking adaptability. In contrast, earlier general-purpose approaches often incurred significant complexity overhead, requiring mechanical movement of optical elements or distinct hardware systems to switch between different tasks. To address this rigidity, emerging multi-wavelength diffractive deep neural networks exploit the spectral dimension to perform parallel classification across multiple databases without mechanical reconfiguration, effectively encoding different inputs into separate wavelength channels. Furthermore, multifunctional processors with learned structured illumination incorporate illumination patterns as trainable parameters, enabling the reuse of diffractive layers to rapidly switch between classification and imaging tasks. However, these flexible multi-layer designs introduce practical challenges regarding fabrication tolerances and opto-mechanical assembly, as they are highly sensitive to layer-to-layer misalignments. Recent advances mitigate this through vaccination strategies that model random displacements during training, yet the core tension remains: maximizing efficiency for specific tasks versus maintaining the versatility required for broader applications, a trade-off that directly influences the physical realizability discussed next. \citep{doi_10_1038_srep21545_ae20cc7d} \citep{doi_10_1515_nanoph_2022_0615_9a3be8d6} \citep{doi_10_1364_ol_591818_8738c4b1} \citep{doi_10_1515_nanoph_2020_0291_c4a9f42d}

Physical implementation bottlenecks ultimately constrain the realization of these diffractive systems, moving beyond the theoretical architectural debate between specialized and flexible designs. Although photonic resonators offer advantageous nonlinear behaviors for optical activation functions, their performance remains fundamentally vulnerable to environmental disturbances such as temperature fluctuations and mechanical vibrations. This sensitivity extends to illumination conditions; real-world deployment demands robust processing of incoherent light, yet most theoretical frameworks rely on idealized quasimonochromatic assumptions that rarely hold outside controlled laboratories. Furthermore, integrating these networks with standard semiconductor processes presents a critical hurdle, as current demonstrations are largely confined to proof-of-concept setups rather than scalable, CMOS-compatible platforms. The challenge compounds when attempting to merge multiple optical functionalities, which often necessitates complex encoding schemes or intricate reconstruction algorithms that undermine practical simplicity. Consequently, while diffractive optical neural networks promise passive energy efficiency and low-latency inference through photon-based computation, achieving fully programmable reconfigurability without succumbing to fixed architecture limitations and alignment sensitivities remains a significant engineering barrier. \citep{doi_10_3390_nano14080697_c372a140} \citep{doi_10_1364_oe_523619_ddeb1b53} \citep{doi_10_1038_s41467_022_28702_0_1b79b43c} \citep{doi_10_1364_optica_583686_6b84680e} \citep{doi_10_1002_lpor_202200348_d54cae18}

In response to the hardware bottlenecks and integration hurdles previously discussed, emerging counter-narratives propose that diffractive neural networks (DONNs) need not rely on complex, multi-layer trainable architectures to be effective. One such approach involves minimalist DONNs, which utilize straightforward optical configurations with efficient computational capabilities, demonstrating robustness and scalability in image classification tasks without the overhead of deep, fully trainable stacks. For instance, a minimalist optical diffractive neural network employs two-level quantized pixel-wise optical encoding to achieve robust image classification without the complexity of deep, fully trainable stacks. \citep{doi_10_1002_lpor_202402303_1f555bd4} Alternatively, end-to-end optoelectronic frameworks reframe the diffractive element not as an active computer but as a fixed passive encoder; in these lensless imaging systems, a designed mask modulates light for downstream algorithmic reconstruction, raising fundamental questions about whether the optical layer itself performs computation or merely conditions the data. In lensless imaging systems, such as those described in advances in mask-modulated techniques, a designed mask acts as a fixed passive encoder that modulates light for downstream algorithmic reconstruction rather than performing active computation. \citep{doi_10_3390_electronics13030617_5b341306} While standard end-to-end differentiable training through the diffraction integral has enabled all-optical generalization to unseen scattering environments, allowing reconstruction through unknown diffusers without retraining, this capability is often predicated on idealized scalar diffraction models. These models frequently neglect material dispersion, fabrication-induced surface roughness, and misalignment errors, creating a significant simulation-to-reality gap. To address this, advanced training schemes now introduce random misalignments as continuous variables during the learning phase, effectively ``vaccinating'' the network to maintain performance despite physical assembly tolerances. Furthermore, attempts to increase flexibility by training different tasks on separate spatial regions of a single layer face hard physical area limitations, typically restricting such multiplexed systems to only two or three simultaneous functions. Collectively, these developments suggest that while ODNNs offer theoretical speed benefits, their practical superiority remains contingent on resolving the tension between minimalist passive encoding and the fragility of complex, idealized active computing. \citep{doi_10_1038_s41377_022_00725_8_57c4c704} \citep{doi_10_1002_lpor_202402303_1f555bd4} \citep{doi_10_48550_arxiv_2411_05748_a91b6f61} \citep{doi_10_3390_electronics13030617_5b341306} \citep{doi_10_1515_nanoph_2020_0291_c4a9f42d}

\section{Bottlenecks and Future Trajectories for Diffractive Neural Networks: Toward Hybrid Electro-Optical Systems and Scalable Optical AI}\label{bottlenecks-and-future-trajectories-for-diffractive-neural-networks-toward-hybrid-electro-optical-systems-and-scalable-optical-ai}

Current progress in transitioning laboratory prototypes to deployable optical AI is hindered by the intrinsic rigidity of static diffractive neural networks, which function as `write-once' hardware incapable of dynamic adaptation without mechanical intervention or complete refabrication. While recent multi-wavelength designs have demonstrated that parallel task processing can alleviate some throughput constraints, these systems remain fundamentally limited in reconfigurability, often requiring separate hardware instances for distinct computational goals. Beyond this lack of flexibility, physical realization introduces severe performance gaps; at shorter wavelengths, even minor fabrication variability in diffractive element sizes can drastically degrade classification accuracy, a challenge compounded by material constraints that limit the number of feasible optical layers. \citep{doi_10_1109_tmag_2023_3281842_3e98eb7f} Furthermore, photonic components such as resonators exhibit acute vulnerability to environmental perturbations, including temperature fluctuations and mechanical vibrations, which destabilize phase relationships and compromise inference reliability. These combined bottlenecks---static functionality, manufacturing sensitivity, and environmental instability---establish a clear boundary for all-optical approaches, necessitating a strategic shift toward hybrid electro-optical co-design to enable the robustness required for real-world deployment. \citep{doi_10_1515_nanoph_2022_0615_9a3be8d6} \citep{doi_10_1002_adpr_202300310_6647aa64} \citep{doi_10_3390_nano14080697_c372a140} \citep{doi_10_1002_nap2_70185_56c8ab30}

A critical simulation-to-reality gap, defined as the performance degradation observed when models trained in idealized numerical environments encounter the noise and aberrations of actual hardware, hinders the transition from theoretical design to physical deployment in diffractive neural networks. While simulations often assume quasimonochromatic illumination with high spatial coherence, real-world applications necessitate the processing of incoherent and partially coherent light, creating a fundamental mismatch in how optical fields evolve through the network layers. \citep{doi_10_1364_oe_523619_ddeb1b53} This discrepancy is exacerbated by fabrication-induced surface roughness, which introduces unmodeled phase errors that significantly degrade inference accuracy compared to digital predictions. Furthermore, compact multi-layer designs face additional constraints from interlayer reflections and interpixel interactions; although these effects may be negligible in low-index materials, they become substantial barriers in high-index systems, limiting the fidelity of analytical models. Addressing these bottlenecks requires moving beyond static design toward fabrication-aware training and hardware-in-the-loop optimization, where online retraining can compensate for physical defects and restore designed functionality. These physical realities underscore the necessity of the hybrid electro-optical strategies discussed next, which aim to mitigate such rigidities by integrating electronic flexibility with optical speed. \citep{doi_10_1038_s41598_025_19638_8_1b2db1c9} \citep{doi_10_1364_oe_523619_ddeb1b53} \citep{doi_10_1109_dac56929_2023_10247715_f78d6713} \citep{doi_10_1364_ol_477605_b9242e3a}

Hybrid electro-optical co-design has emerged as a dominant strategy to address the reconfigurability and nonlinearity bottlenecks identified in static hardware. This architecture integrates static diffractive layers, which leverage the inherent parallelism and low latency of light propagation for high-throughput feature extraction, with reconfigurable electronic layers responsible for nonlinear activation and logic operations. By offloading these computationally intensive nonlinear tasks to electronics, such systems mitigate the rigidity of all-optical hardware while preserving significant energy efficiency gains. Representative implementations include programmable diffractive deep neural networks based on digital-coding metasurface arrays, which retain the flexibility of electronic control to overcome the fixed network architectures of passive devices. Programmable diffractive deep neural networks have been realized using digital-coding metasurface arrays, which leverage electronic control to maintain flexibility and overcome the static nature of passive optical devices. \citep{doi_10_1038_s41928_022_00719_9_a2169947} Furthermore, research demonstrates that appending convolutional neural networks to the terminal of all-optical multilayered networks allows for joint training, where electronic components introduce necessary nonlinear functions like ReLU and Sigmoid to enhance overall system capability. While purely optical approaches struggle with arbitrary modulation and precise control, this hybrid pathway effectively bridges the gap between optical speed and electronic adaptability, setting the stage for further dynamic control mechanisms via programmable materials. \citep{doi_10_3390_nano14221812_4d7d91b1} \citep{doi_10_1038_s41928_022_00719_9_a2169947} \citep{doi_10_1126_sciadv_adu0904_786138ff}

To overcome static hardware constraints, recent approaches incorporate dynamic adaptability via programmable materials and online learning. Programmable metasurfaces, defined as dynamic optical surfaces capable of altering phase profiles via external control signals, address the rigidity of passive diffractive layers. For instance, magneto-optical diffractive deep neural networks utilize the Faraday effect, where the magnetization direction of magnetic domains modulates light phase without requiring complex microfabrication, thus allowing for online learning techniques that dynamically reconfigure the network to correct for alignment or fabrication errors. Complementing this linear reconfigurability, the integration of tunable nonlinear activation functions overcomes the bottleneck of fixed optical nonlinearity. By employing photochromic materials like WO3 films, these systems leverage UV excitation to induce intervalent electron transfer, reversibly altering the film's complex amplitude distribution to emulate adaptive neuron-like operations. Adaptive diffractive deep neural networks have utilized WO\(_{3}\) films to enable tunable nonlinear activation functions, where UV excitation induces intervalent electron transfer to reversibly alter the film's complex amplitude distribution. \citep{doi_10_1002_lpor_202500558_72fd7a4f} While such physics-aware training frameworks facilitate the deployment of computer-trained models onto reconfigurable hardware, these designs primarily serve as a viable near-term trajectory; they leverage optical speed for inference while retaining the electronic flexibility necessary for rapid updates and error correction, rather than achieving fully autonomous all-optical programming. \citep{s2_44f5e5b7573063092084a20caab78ea8a4e2407e_dd441649} \citep{doi_10_1002_lpor_202500558_72fd7a4f} \citep{doi_10_1002_lpor_202200348_d54cae18}

Programmable materials address reconfigurability, yet the operational bandwidth of diffractive systems remains constrained by wavelength-dependent phase delays, necessitating advanced design strategies to support hyperspectral applications. Multi-element metasurface systems overcome these limitations by harnessing illumination angle-dependent phase delays across local units, effectively decoupling angular degrees of freedom to enable robust holographic imaging without relying on complex intrinsic angular dispersion. Researchers demonstrated a multiwavelength achromatic deflector in the visible spectrum using a single-layer freeform metasurface to achieve broadband operation without relying on complex intrinsic angular dispersion. \citep{doi_10_1021_acs_nanolett_4c02995_4955c826} Complementing this structural approach, data-driven inverse design frameworks utilize differentiable neural simulators to map metasurface geometries directly to spectral responses, allowing for the joint end-to-end optimization of optical encoders and reconstruction networks. This methodology has demonstrated significant improvements in spectral fidelity and noise robustness compared to traditional unoptimized filters. The critical advantage of expanding bandwidth is further evidenced in scattering environments, where broadband single-pixel diffractive networks achieved blind testing accuracies of approximately 88.53\% through unknown random diffusers, clearly surpassing the performance of monochrome designs that suffered a roughly 10\% accuracy deficit. These findings underscore that wavelength multiplexing offers significant advantages for information processing capacity, while realizing these designs at scale requires embedding physical fabrication constraints directly into the training loop to ensure experimental viability. \citep{doi_10_2528_pier25052305_11655dbe} \citep{s2_76a83ea3baed75ad88b7f0782eb9a694a330c853_cee16cac} \citep{doi_10_1038_s41377_023_01116_3_aba67a18}

Bridging the persistent simulation-to-reality gap requires optimization protocols to evolve from idealized numerical emulation to physics-aware training that embeds hardware constraints directly into the learning loop. A primary challenge arises because directly mapping discrete analog device responses to diffractive optical neural network (DONN) systems breaks the gradient chain required for backpropagation, as optical components exhibit non-uniform, non-monotonic behaviors unlike stable digital computation. Physics-aware differentiable discrete codesign addresses this by utilizing techniques such as Gumbel-Softmax to enable fully differentiable mappings that respect physical periodicity and device-specific ranges, thereby narrowing the discrepancy between training and deployment. \citep{doi_10_1145_3508352_3549378_13c53996} Beyond discrete mapping, fabrication-induced phase errors and wavefront aberrations remain critical bottlenecks that degrade fidelity in dielectric metasurfaces. Aberration-adaptive learning mitigates these issues by incorporating generated wavefront distortions, such as those from Zernike polynomials, into every training batch. \citep{s2_8560e3428b029604da9bd494e485939165b3777e_bf7326ee} The configuration forces the model to adapt to a wide range of unknown aberrations rather than optimizing for a single ideal state, experimentally stabilizing classification accuracy against variations caused by material shrinkage or thermal expansion. By co-optimizing the optical design with these inevitable physical imperfections, researchers can reduce the precision requirements for fabrication while ensuring robust performance, a necessary step toward the hybrid electro-optical systems proposed as the viable path forward for scalable optical AI. \citep{doi_10_1145_3508352_3549378_13c53996} \citep{s2_8560e3428b029604da9bd494e485939165b3777e_bf7326ee} \citep{doi_10_1063_5_0141881_93942631}

Even with physics-aware optimization protocols embedding constraints directly into the training loop, ODNN research lacks standardized benchmarking datasets necessary to validate progress beyond narrow 2D spatial imaging tasks. Current high-fidelity resources, while rigorous, remain restricted to fixed cellular specimens and widefield-based architectures, failing to encapsulate the spatiotemporal dynamics or volumetric deconvolution challenges required for broader deployment. This data gap is compounded by severe computational bottlenecks in online in-situ training, where compensating for system errors in high-dimensional datasets like MNIST demands processing times that exceed the stability windows of experimental setups affected by temperature and vibration. Consequently, researchers are often forced to rely on offline error compensation rather than dynamic adaptation, highlighting a critical barrier to maturity. Addressing these limitations through unified evaluation protocols and efficient training strategies is essential to transition diffractive neural networks from static laboratory prototypes toward the hybrid electro-optical systems envisioned for scalable optical AI. \citep{doi_10_1038_s41597_026_07467_x_ae097478} \citep{s2_dac795899d5e98d409739c3bcfe096491d6b7427_5b5c8d05}

Near-term engineering strategies prioritize transfer learning and multifunctional processor designs to bridge the gap between static hardware limitations and the need for dynamic functionality. Transfer learning addresses the rigidity of fixed physical structures by incorporating partially reconfigurable optical units that allow pre-trained layers to adapt to new tasks without full re-fabrication; for instance, numerical simulations of such architectures have demonstrated multi-task processing with classification accuracies reaching 98.8\% on MNIST datasets. These designs validate that minimalist optical configurations can be effectively deployed for specific inference applications, including holographic data storage and image classification, leveraging their inherent computational efficiency despite limited reconfigurability. While early all-optical diffractive deep neural networks successfully demonstrated ultra-low latency and high parallelism using free-space light interconnection, they remain constrained by fixed hardware architectures that prevent post-fabrication updates. Consequently, the realization of fully programmable all-optical networks remains experimentally unachieved, primarily due to the physical challenges of integrating dynamic, low-loss phase modulation into dielectric metasurfaces at scale. Although all-dielectric metasurface-based designs offer a path toward miniaturization and higher spatial neuron density by reducing component sizes below the wavelength scale, they currently serve as proof-of-concept steps rather than complete solutions. Thus, the immediate trajectory favors hybrid electro-optical co-design to inject dynamic functionality, reserving the goal of entirely programmable all-optical systems for long-term material and fabrication breakthroughs. \citep{doi_10_1117_12_3109287_f4650a31} \citep{doi_10_3390_photonics11020145_68f5de58} \citep{doi_10_1364_josab_497148_b087424b} \citep{doi_10_1117_12_2688132_926a2697}

All-optical diffractive classifiers offer distinct advantages in real-time and energy-efficient inference compared to conventional convolutional neural networks, making them attractive for applications where speed is critical, even though they currently lag in raw accuracy. While diffractive deep neural networks (D2NNs) provide ultra-low latency and low power consumption through highly parallel optical computing, their practical deployment faces significant bottlenecks, including fabrication-induced phase errors, alignment challenges across multiple diffractive planes, and limited operational bandwidth. Furthermore, current research indicates that training different tasks on separate regions of a static diffractive layer is constrained by area limitations, typically restricting a single device to only two or three fixed functions. These rigidity issues highlight why a purely optical approach remains insufficient for scalable multi-function AI, reinforcing the chapter's thesis that hybrid electro-optical architectures represent the most viable trajectory. By leveraging optical speed for linear sub-tasks while relying on electronic flexibility for error correction and reconfigurability, such hybrid systems can overcome the inherent trade-offs between inference speed and functional adaptability. \citep{doi_10_1364_ol_593231_ab93febc} \citep{doi_10_1364_oe_598903_731aaac1} \citep{doi_10_48550_arxiv_2411_05748_a91b6f61}

\section{Conclusion}\label{conclusion}

This review synthesizes the current state of Optical Diffractive Neural Networks (ODNNs), addressing the central inquiry regarding their viability as a computational paradigm distinct from conventional electronics. The evidence presented across the physical, training, and hardware domains establishes that ODNNs function not merely as optical accelerators but as unique physical computers where the laws of diffraction dictate computational expressivity. By replacing electronic weight matrices with spatially varying phase profiles, these systems leverage the inherent parallelism of light to perform linear transformations at the speed of light with minimal energy consumption. This architecture offers a compelling solution to the von Neumann bottleneck, enabling ultra-low latency inference for specific tasks such as real-time image classification, all-optical encryption, and compressive sensing. However, the body of literature clarifies that this advantage is currently domain-specific; ODNNs excel in throughput and energy efficiency for linear or pre-defined operations but do not yet constitute a general-purpose replacement for electronic processors due to fundamental physical constraints.

Tension between idealized mathematical models and physical reality constitutes the primary limitation preventing widespread deployment. A critical barrier is the simulation-to-reality gap, where performance predicted by scalar diffraction models degrades in physical prototypes due to fabrication tolerances, surface roughness, and material dispersion. While techniques such as Hardware-in-the-Loop (HIL) optimization and ``vaccination'' against misalignment have shown promise in bridging this divide, they underscore the fragility of current designs. Furthermore, the inherent linearity of passive diffractive layers imposes a nonlinearity bottleneck, restricting the network's ability to solve complex problems without external intervention. This limitation, combined with the static nature of most fabricated metasurfaces and 3D-printed elements, results in ``write-once'' hardware that lacks the reconfigurability essential for dynamic computing environments. Consequently, purely all-optical systems remain vulnerable to environmental perturbations, such as temperature fluctuations and mechanical vibrations, which can destabilize phase relationships and compromise inference reliability.

Looking forward, the trajectory for ODNN research appears to be shifting away from the pursuit of entirely all-optical general-purpose computers toward hybrid electro-optical co-design. The most viable path for scalable optical AI involves leveraging the speed and parallelism of diffractive layers for linear feature extraction while offloading nonlinear activation functions and adaptive control to electronic components. This hybrid approach mitigates the rigidity of static hardware and allows for error correction and dynamic task switching that purely optical systems cannot currently support. Additionally, future progress depends on the adoption of physics-aware training paradigms that embed fabrication constraints and device non-idealities directly into the loss landscape, ensuring that designed networks are robust to manufacturing variances. To validate these advancements, ODNN research must also establish standardized benchmarking protocols that move beyond narrow 2D imaging tasks to evaluate spatiotemporal dynamics and volumetric processing under realistic conditions. In conclusion, while ODNNs have matured from theoretical concepts to functional laboratory prototypes demonstrating unique physical advantages, their transition to robust, deployable technologies will rely on integrating optical speed with electronic flexibility, thereby creating resilient systems capable of operating effectively outside controlled laboratory environments.

\bibliographystyle{unsrtnat}
\bibliography{references}
\end{document}
